% sec12_5.tex — Section 12.5: Method of Characteristics for a Vibrating String of Fixed Length
\section{Method of Characteristics for a Vibrating String of Fixed Length}
\label{sec:char-fixed-string}

In Chapter~2 we solved for the vibration of a finite string satisfying
\boxed{\begin{equation}
\label{eq:12.5.1}
\text{PDE:}\quad \pdd{u}{t} = c^2\pdd{u}{x}
\end{equation}}
\boxed{\begin{equation}
\label{eq:12.5.2}
\text{BC:}\quad u(0,t) = 0, \qquad u(L,t) = 0
\end{equation}}
\boxed{\begin{equation}
\label{eq:12.5.3}
\text{IC:}\quad u(x,0) = f(x), \qquad \pd{u}{t}(x,0) = g(x),
\end{equation}}
using Fourier series methods.  We can obtain an equivalent, but in some ways more useful, result by using the general solution of the one-dimensional wave equation:
\begin{equation}
\label{eq:12.5.4}
u(x,t) = F(x - ct) + G(x + ct).
\end{equation}
The initial conditions are prescribed only for $0 < x < L$, and hence the formulas for $F(x)$ and $G(x)$ previously obtained are valid only for $0 < x < L$:
\begin{align}
\label{eq:12.5.5}
F(x) &= \frac{1}{2}f(x) - \frac{1}{2c}\int_0^x g(\bar{x})\,d\bar{x} \\
\label{eq:12.5.6}
G(x) &= \frac{1}{2}f(x) + \frac{1}{2c}\int_0^x g(\bar{x})\,d\bar{x}.
\end{align}

\begin{figure}[H]
\centering
\includegraphics[width=0.8\textwidth]{figures/ch12/fig_12_5_1.png}
\caption{Characteristics.}
\label{fig:12.5.1}
\end{figure}

If $0 < x - ct < L$ and $0 < x + ct < L$, as shown shaded in Fig.~12.5.1, then d'Alembert's solution is valid:
\begin{equation}
\label{eq:12.5.7}
u(x,t) = \frac{f(x-ct) + f(x+ct)}{2} + \frac{1}{2c}\int_{x-ct}^{x+ct} g(\bar{x})\,d\bar{x}.
\end{equation}
In this region the string does not know that either boundary exists; the information that there is a boundary propagates at velocity $c$ from $x = 0$ and $x = L$.

If one's position and time is such that signals from the boundary have already arrived, then modifications in \eqref{eq:12.5.7} must be made.  The boundary condition at $x = 0$ implies that
\begin{equation}
\label{eq:12.5.8}
0 = F(-ct) + G(ct) \qquad \text{for } t > 0,
\end{equation}
while at $x = L$ we have
\begin{equation}
\label{eq:12.5.9}
0 = F(L - ct) + G(L + ct) \qquad \text{for } t > 0.
\end{equation}
These, in turn, imply reflections and multiple reflections, as illustrated in Fig.~12.5.2.

\begin{figure}[H]
\centering
\includegraphics[width=0.8\textwidth]{figures/ch12/fig_12_5_2.png}
\caption{Multiply reflected characteristics.}
\label{fig:12.5.2}
\end{figure}

Alternatively, a solution on an infinite domain without boundaries can be considered that is odd around $x = 0$ and odd around $x = L$, as sketched in Fig.~12.5.3.  In this way, the zero condition at both $x = 0$ and $x = L$ will be satisfied.  We note that $u(x,t)$ is periodic with period $2L$, since periodic functions that are odd around $x = 0$ are automatically odd around $x = L$.  Thus, the simplest way to obtain the solution is to \textbf{extend the initial conditions as odd functions (around $x = 0$) which are periodic (with period $2L$)}.  With these odd periodic initial conditions, the method of characteristics can be utilized as well as d'Alembert's solution \eqref{eq:12.5.7}.

\begin{figure}[H]
\centering
\includegraphics[width=0.8\textwidth]{figures/ch12/fig_12_5_3.png}
\caption{Odd periodic extension.}
\label{fig:12.5.3}
\end{figure}

\textbf{Example.}  Suppose that a string is initially at rest with prescribed initial conditions $u(x,0) = f(x)$.  The string is fixed at $x = 0$ and $x = L$.  Instead of using Fourier series methods, we extend the initial conditions as odd functions around $x = 0$ and $x = L$.  Equivalently, we introduce the \emph{odd periodic extension}.  (The odd periodic extension is also used in the Fourier series solution.)  Since the string is initially at rest, $g(x) = 0$; the odd periodic extension is $g(x) = 0$ for all $x$.  Thus, the solution of the one-dimensional wave equation is the sum of two simple waves:
\[
u(x,t) = \frac{1}{2}[f_{\text{ext}}(x - ct) + f_{\text{ext}}(x + ct)],
\]
where $f_{\text{ext}}(x)$ is the odd periodic extension of the given initial position.  This solution is much simpler than the summation of the first 100 terms of its Fourier sine series.

\textbf{Separation of variables.}  By separation of variables, the solution of the wave equation with fixed boundary conditions ($u(0,t) = 0$ and $u(L,t) = 0$) satisfying the initial conditions $u(x,0) = f(x)$ and $\frac{\partial u}{\partial t} = 0$ is
\[
u(x,t) = \sum_{n=1}^{\infty} a_n \sin\frac{n\pi x}{L}\cos\frac{n\pi ct}{L},
\]
where the given initial conditions are only specified for $[0,L]$.  To be precise, the infinite Fourier series $f_{\text{ext}}(x)$ equals the odd periodic extension (with period $2L$) of $f(x)$: $f_{\text{ext}}(x) = \sum_{n=1}^{\infty} a_n \sin\frac{n\pi x}{L}$.  Using $\sin\frac{n\pi x}{L}\cos\frac{n\pi ct}{L} = \frac{1}{2}\sin\frac{n\pi(x+ct)}{L} + \frac{1}{2}\sin\frac{n\pi(x-ct)}{L}$, we obtain
\[
u(x,t) = \frac{1}{2}f_{\text{ext}}(x + ct) + \frac{1}{2}f_{\text{ext}}(x - ct),
\]
which is the same result obtained by the method of characteristics.

\begin{exercises}{12.5}

\item \label{ex:12.5.1}
Consider
\[
\pdd{u}{t} = c^2\pdd{u}{x}
\]
\begin{align*}
u(x,0) &= f(x) \\
\pd{u}{t}(x,0) &= g(x)
\end{align*} \quad $0 < x < L$
\begin{align*}
u(0,t) &= 0 \\
u(L,t) &= 0.
\end{align*}
\begin{enumerate}
\item[(a)] Obtain the solution by Fourier series techniques.
\item[(b)]\starred\ If $g(x) = 0$, show that part (a) is equivalent to the results of Chapter~12.
\item[(c)] If $f(x) = 0$, show that part (a) is equivalent to the results of Chapter~12.
\end{enumerate}

\item \label{ex:12.5.2}
Solve using the method of characteristics:
\[
\pdd{u}{t} = c^2\pdd{u}{x}
\]
\begin{align*}
u(x,0) &= 0 \qquad & u(0,t) &= h(t) \\
\pd{u}{t}(x,0) &= 0 \qquad & u(L,t) &= 0.
\end{align*}

\item \label{ex:12.5.3}
Consider
\[
\pdd{u}{t} = c^2\pdd{u}{x} \qquad 0 < x < 10
\]
\[
u(x,0) = f(x) = \begin{cases} 1 & 4 < x < 5 \\ 0 & \text{otherwise} \end{cases} \qquad u(0,t) = 0
\]
\[
\pd{u}{t}(x,0) = g(x) = 0 \qquad \pd{u}{x}(L,t) = 0.
\]
\begin{enumerate}
\item[(a)] Sketch the solution using the method of characteristics.
\item[(b)] Obtain the solution using Fourier-series-type techniques.
\item[(c)] Obtain the solution by converting to an equivalent problem on an infinite domain.
\end{enumerate}

\item \label{ex:12.5.4}
How should initial conditions be extended if $\partial u/\partial x(0,t) = 0$ and $u(L,t) = 0$?

\end{exercises}
